\section{Conclusion}
\label{sec:conclusion}

We tested whether moving the published 4-factor MCDM lending
allocator \cite{solovev2026vault} from hourly polling to per-block
event-time evaluation, paired with HFT-inspired gas-aware switching
rules, recovers the missing edge that the hourly version's H$_{0}$
result \cite{solovev2026lob} exposed.  On a $\NBlocksFull$-block
panel covering \DataStart--\DataEnd\ across \NProtocols\ Ethereum
USDC lending venues (\ProtocolList), the gas-aware threshold
allocator (T1) outperforms the published hourly MCDM-EMA baseline by
$+20$\,bp net APY ($\TOneAPY$ vs $\EMAAPY$) while issuing
$\TOneNRebal$ rebalances over four months --- roughly twenty times
the published baseline's two-in-four-months count, at one-third the
gas cost per switch of the unrefined greedy benchmark.

Against passive buy-and-hold per protocol, T1 earns
\$4{,}275 above Aave V3 hold ($+42.7$\,bp, $H_{1}^{\text{aux}}$
$\Delta\textsc{Sharpe} = \HOneAuxDelta$, $p = \HOneAuxPvalue$),
\$4{,}039 above Morpho Blue hold ($+40.4$\,bp), and trails Euler V2
hold by \$817 ($-8.2$\,bp).  The Euler-side shortfall is the natural
journal-extension hook: closing it requires the F1 lead-rate panel
(Maker DSR, Curve 3pool) and a cross-protocol lead signal that
specifically anticipates Euler-leads-Aave concentration windows.
The architecturally identical Cox proportional-hazards policy (T3) is
trained on the panel with the F3 fragmentation family alone (C-index
$0.628$), and the H$_{1}^{c}$ contrast (T3 vs T2) is reported in
\Cref{tab:h1-significance}; absent the F1 lead-rate covariate T3 does
not yet separate from T2.

The contribution stands at three layers.  At the methodological
layer, the event-time pivot is empirically validated: per-block
gas-aware decisions admit two to three orders of magnitude more
profitable rebalance opportunities than hourly polling, and the same
4-factor MCDM aggregation gains directional edge purely from the
resolution change.  At the infrastructure layer, the seven-way
event-stream fetcher harness, the per-block replay engine, and the
zero-drift production agent (Windows directory junction between
research and live decision modules) are open-source and reproducible
end-to-end (artifact: \texttt{submission\_<git-sha>.zip} with
deterministic sha256 sidecar and per-file manifest).  At the
empirical layer, the binding finding --- the event-time MCDM
allocator significantly beats passive Aave hold under a paired
monthly Sharpe bootstrap --- complements the falsified hourly
LOB-to-DeFi transfer of \cite{solovev2026lob} and the equilibrium
impossibility result of \cite{halpern2024fair} by demonstrating that
the residual exploitable edge lives in the inter-block
microstructure, not in the hourly forecast.

Three concrete next steps follow.  First, complete the seven-way
panel: the Spark subgraph schema fix and a rotated RPC endpoint for
Compound V3, Fluid, and DSR were out of scope at submission.  Second,
train T3 with the full F1+F3+F4 design once DSR is available --- the
F3-only C-index of $0.628$ is a lower bound on what the hazard model
can attain.  Third, run the production agent against Sepolia for the
$\ge 10$-rebalance acceptance gate documented in
\texttt{agent/RUNBOOK.md}; the off-chain agent is feature-complete
and the Flashbots private-mempool dispatch path is dry-run-verified
end-to-end.  Together these close the loop from research to
deployment on a methodology --- per-block, gas-aware,
literature-grounded multi-protocol allocation --- whose hourly
predecessor could not have been falsified or vindicated without the
resolution change.
